% Section: EMF chain map. \input by report_2.tex. Detail in app:emf.
\section{EMF chain}\label{sec:emf}

\paragraph{One problem, followed end to end.}
Take the smallest instance, \texttt{emf/instances/1-path.exam}: four exams whose conflicts form a path,
two batches.

\VerbatimInput[fontsize=\footnotesize, frame=topline, label=\texttt{1-path.exam}]{../../emf/instances/1-path.exam}

Running \texttt{cd emf \&\& ./build.sh} turns that text into the files in \texttt{emf/out/path/}. What
happens to it, in order:

\begin{enumerate}[nosep, leftmargin=*]
	\item \texttt{1-path.exam} (text).
	\item Xtext: the parser from \texttt{Exam.xtext} (text~$\to$~model) yields \texttt{problem.xmi} and
	      \texttt{solverconfig.xmi}, which conform to \texttt{problem.ecore} and \texttt{solverconfig.ecore}.
	\item OCL validation, then ATL \texttt{problem2conflictgraph.atl} (model~$\to$~model) yields
	      \texttt{conflictgraph.xmi}, which conforms to \texttt{conflictgraph.ecore}.
	\item OCL validation, then Acceleo \texttt{main.mtl} (model~$\to$~text) yields
	      \texttt{solve\_path.py}.
	\item Python at run time: pyqubo builds the QUBO, then QAOA or annealing samples it, writing
	      \texttt{schedule\_path.xmi} and a printed report; the schedule conforms to \texttt{schedule.ecore}.
\end{enumerate}

The Java that glues the first three arrows is one class,
\texttt{ChainApplication.java} in the chain bundle
(\path{emf/bundles/pt.up.feup.examschedule.chain/}). Its \texttt{run} method is the chain above: parse,
validate, save, ATL, save, validate, Acceleo. The grammar, ATL module, Acceleo template and runtime are
expanded in \cref{app:emf}.

\paragraph{Three metalevels.}
Each \texttt{.xmi} file is a \emph{model}. It is only meaningful because it \emph{conforms to} a
\emph{metamodel}, an \texttt{.ecore} file, which in turn is written in Ecore, EMF's metamodelling language;
XMI is how EMF serialises those models.

\begin{center}
	\small
	\begin{tabular}{@{}lll@{}}
		\toprule
		Level & In this repo & Says \\
		\midrule
		Meta-metamodel & Ecore itself & there are classes, attributes, references \\
		Metamodel & \texttt{problem.ecore} & a Problem has exams; an Exam has students \\
		Model & \texttt{emf/out/path/problem.xmi} & Calculus is sat by Ana and Rui \\
		\bottomrule
	\end{tabular}
\end{center}

The OCL constraints (\texttt{uniqueExamNames}, \texttt{atLeastOneBatch},\ldots) live \emph{inside} the
metamodel as annotations, so they travel with it and every tool that loads \texttt{problem.ecore} can check
them (\cref{app:metamodels}).

\paragraph{Three kinds of arrow.}
Model-driven engineering names transformations by what goes in and out. The chain has one of each:

\begin{itemize}[nosep]
	\item \textbf{Text~$\to$~model}: Xtext. The grammar does not invent its own classes; its rules say
	      \texttt{returns problem::Problem}, so the parser builds instances of the hand-written metamodel.
	\item \textbf{Model~$\to$~model}: ATL. Rules map a Problem to a Conflict Graph and add an edge for each
	      pair of exams sharing a student, the paper's Algorithm~1 \cite[§5.2, Algorithm 1]{Paper}.
	\item \textbf{Model~$\to$~text}: Acceleo. A template walks the Conflict Graph and prints Python.
\end{itemize}

The last model the chain itself produces is the Conflict Graph. The QUBO never becomes an \texttt{.xmi}:
pyqubo builds it in memory when the generated script runs (\texttt{emf/docs/adr/0007}), the hand-off chosen
in \cref{sec:strategy}.
